\subsection{Resumo}

Este primeiro capítulo argumentou a favor de um repensar da filosofia estatística e científica popular. Em vez de escolher entre várias ferramentas de caixa-preta para testar hipóteses nulas, devemos aprender a construir e analisar múltiplos modelos não nulos de fenômenos naturais. Para apoiar este objetivo, o capítulo introduziu a inferência bayesiana, a comparação de modelos, os modelos multinível e os modelos causais gráficos. O restante do livro está organizado em quatro partes interdependentes.

\begin{enumerate}
  \item Os Capítulos 2 e 3 são fundamentais. Eles introduzem a inferência bayesiana e as ferramentas básicas para realizar cálculos bayesianos. Eles avançam bem devagar e enfatizam uma interpretação puramente lógica da teoria das probabilidades.
  \item Os próximos quatro capítulos, do 4 ao 9, constroem a regressão linear múltipla como uma ferramenta bayesiana. Esta ferramenta apoia a inferência causal, mas apenas quando analisamos modelos causais separados que nos ajudam a determinar quais variáveis incluir. Estes capítulos enfatizam a plotagem de resultados em vez de tentar interpretar estimativas de parâmetros individuais. Problemas de complexidade do modelo — \textsc{overfitting} — também aparecem com destaque. Portanto, você também terá uma introdução à teoria da informação e à comparação formal de modelos no Capítulo 7.
  \item A terceira parte do livro, Capítulos 9 a 12, apresenta modelos lineares generalizados de vários tipos. O Capítulo 9 introduz o método de Monte Carlo via cadeias de Markov (\textit{Markov chain Monte Carlo}), usado para ajustar os modelos nos capítulos posteriores. O Capítulo 10 introduz a entropia máxima como um procedimento explícito para nos ajudar a projetar e interpretar esses modelos. Em seguida, os Capítulos 11 e 12 detalham os próprios modelos.
  \item A última parte, Capítulos 13 a 16, aborda modelos multinível, bem como modelos especializados que lidam com erros de medição, dados ausentes e correlação espacial. Este material é bastante avançado, mas prossegue da mesma maneira mecanicista que o material anterior. O Capítulo 16 afasta-se do restante do livro ao implantar modelos que não são do tipo linear generalizado, mas sim modelos teóricos expressos diretamente como modelos estatísticos.
\end{enumerate}

O capítulo final, Capítulo 17, retorna a algumas das questões levantadas neste primeiro.

No final de cada capítulo, há problemas práticos que variam do fácil ao difícil. Esses problemas ajudam você a testar sua compreensão. Os mais difíceis expandem o material, introduzindo novos exemplos e obstáculos.